\documentclass[10pt,a4paper]{article}
\usepackage[T1]{fontenc}
\usepackage[utf8]{inputenc}
\usepackage[top=1.6cm,bottom=1.6cm,left=1.8cm,right=1.8cm]{geometry}
\usepackage{titlesec}
\usepackage{enumitem}
\usepackage[hidelinks]{hyperref}
\usepackage{textcomp}

\pagestyle{empty}
\setlength{\parindent}{0pt}

\titleformat{\section}{\large}{}{0em}{}[\vspace{-0.9em}\noindent\rule{\textwidth}{0.4pt}]
\titlespacing*{\section}{0pt}{1.0em}{0.5em}

\setlist[itemize]{leftmargin=1.5em,itemsep=0.1em,topsep=0.3em,parsep=0pt,label=$\bullet$}

% \entry{Title}{Right}{Subtitle}{Right subtitle}
\newcommand{\entry}[4]{%
  \textbf{#1}\hfill\textbf{#2}\\
  {\small\textit{#3}}\hfill{\small\textit{#4}}\par}

\begin{document}

\begin{center}
  {\fontsize{24}{28}\selectfont Mitnala Shiva Jyothiradithya}\\[0.5em]
  B.Tech in Mathematics and Computing $|$ Indian Institute of Science, Bengaluru\\[0.2em]
  {\small \href{mailto:jyothiradit1@iisc.ac.in}{jyothiradit1@iisc.ac.in} \quad$|$\quad
  \href{https://github.com/msja1729}{github.com/msja1729} \quad$|$\quad Bengaluru, India}
\end{center}

\section*{Education}
\entry{Indian Institute of Science, Bengaluru}{2024 -- 2028}{B.Tech in Mathematics and Computing}{CGPA: 8.9 / 10 (through Semester 4)}
\entry{Sri Chaitanya Junior College}{2024}{Class XII -- Board of Intermediate Education, Andhra Pradesh}{Percentage: 99.1}
\entry{Montessori Indus EM Primary School}{2022}{Class X -- CBSE}{Percentage: 98.4}

\section*{Experience}
\entry{Emergence AI}{May 2026 -- Jul 2026}{Research Intern}{}
\begin{itemize}
  \item Architected \textbf{BulletproofTemplateS}, a context-aware web templating engine built from scratch in Lean, providing mathematically rigorous guarantees against Cross-Site Scripting (XSS) vulnerabilities.
  \item Engineered and formally verified a custom recursive-descent parser and lexer that processes HTML, CSS, JavaScript, JSON, and Jinja-style logic into an Abstract Syntax Tree (AST).
  \item Developed a context-aware escaping engine that structurally analyzes the AST to automatically apply correct sanitization rules across HTML, JS, CSS, and JSON contexts.
  \item Verified overall correctness, security, and stability across all possible inputs using machine-checked proofs, guaranteeing safety and liveness properties across Web and LLM environments.
\end{itemize}

\section*{Projects}
\entry{SMRITI: Memory Firewall for AI Agents}{AEGIS 2026, IISc Bengaluru}{Team Lead (4 members) $\cdot$ Python, LangGraph, PostgreSQL + pgvector, NeMo Guardrails, Presidio, Lean 4}{}
\begin{itemize}
  \item Built a provenance-based security layer that protects LLM agents from \textbf{memory-poisoning attacks}, where stored malicious text triggers harmful actions (e.g., fraudulent payments) weeks later.
  \item Designed three checkpoints: a \textbf{write gate} (source-based trust labels, quarantine of instruction-like untrusted text, PII redaction), a \textbf{read gate} that admits only trusted memory into agent instructions, and an \textbf{action check} requiring human approval for risky actions.
  \item Created \textbf{MemPoison-Bench} (60 attacks); cut attack success from \textbf{48\% to 0/60} vs.\ 23\% for NVIDIA NeMo Guardrails, with \textbf{100\% clean-task completion} and no measurable added latency.
  \item Built a counterfactual \textbf{Immune Auditor} that traces a harmful action to its source memory and rolls back it and all derived memory, with \textbf{0 false rollbacks}.
  \item \textbf{Formally verified} core guarantees in Lean 4 (no trust escalation; summaries inherit lowest input trust), backed by 460+ unit tests and differential testing against the proofs.
\end{itemize}

\entry{Fine-Tuning of a Liquid Foundation Model}{May 2026 -- Jun 2026}{ML Researcher}{}
\begin{itemize}
  \item Engineered a custom reinforcement learning pipeline, integrating a GRPO trainer with self-designed reward functions for logical correctness and strict format adherence.
  \item Applied the pipeline to fine-tune a Liquid Foundation Model, substantially improving its mathematical reasoning and problem-solving performance on the GSM8k benchmark.
  \item Evaluated the model across standard industry benchmarks, demonstrating measurable gains in mathematical reasoning and output formatting.
\end{itemize}

\entry{XAI and Gradient Boosted Tree Models}{Mar 2026 -- Apr 2026}{ML Researcher}{}
\begin{itemize}
  \item Replicated foundational XGBoost benchmarks to validate core optimizations, including second-order Taylor approximations and sparsity-aware split finding across diverse predictive tasks.
  \item Benchmarked exact greedy versus approximate methods on the HIGGS dataset (11M events) for competitive AUC, and analyzed Learning-to-Rank algorithms on MQ2008, confirming a 1\% NDCG@10 improvement with 2x faster training.
  \item Evaluated sparsity handling using the Adult dataset and a 10k subset of the Allstate dataset, demonstrating consistent precision and significant reductions in training time.
  \item Analyzed Disk I/O versus compute trade-offs, measuring linear loading scalability ($\sim$1.6M rows/s) and profiling block compression overheads in RAM-heavy environments.
\end{itemize}

\section*{Talks \& Presentations}
\entry{PageRank Algorithm and the Ergodic Theorem}{}{Technical Talk, Indian Institute of Science}{}
\begin{itemize}
  \item Delivered a technical presentation on the PageRank algorithm formulated as a Markov chain.
  \item Explained the role of stochastic matrices, stationary distributions, and eigenvectors in ranking.
  \item Discussed convergence guarantees using the ergodic theorem and irreducibility assumptions, and analyzed the effect of the damping factor on convergence speed and ranking stability.
\end{itemize}

\section*{Achievements}
\begin{itemize}
  \item Secured State 2\textsuperscript{nd} Rank in the State Board Intermediate Examinations.
  \item Secured All India Rank (AIR) 1529 and EWS Rank 180 in JEE Main.
  \item Secured All India Rank (AIR) 1903 and EWS Rank 175 in JEE Advanced.
  \item Awarded 2\textsuperscript{nd} Place in the NPCI Hackathon conducted at the Indian Institute of Science (IISc).
  \item Won 3\textsuperscript{rd} Place as team lead in the AEGIS 2026 Hackathon conducted by Infosys and NVIDIA.
  \item Current recipient of the KVB Scholarship Program (Rs.\ 1,00,000/year), awarded by Karur Vysya Bank.
\end{itemize}

\section*{Technical Skills}
\begin{itemize}
  \item \textbf{Programming:} Python (NumPy, SciPy, Matplotlib, PyTorch), C++ (STL), Lean
  \item \textbf{Mathematics:} Linear Algebra, Probability \& Statistics, Optimization, Numerical Methods, Real Analysis
  \item \textbf{Algorithms:} Sorting, Trees, Graph Algorithms, Dynamic Programming, Gradient Descent, XGBoost
  \item \textbf{Tools \& Frameworks:} FastAPI, SQL, GRPO / RL Pipelines, Vision-Language Models, LangGraph, pgvector
\end{itemize}

\section*{Relevant Coursework}
Real Analysis $\bullet$ Linear Algebra I \& II $\bullet$ Discrete Mathematics $\bullet$ Probability and Statistics $\bullet$ Numerical Methods $\bullet$ Data Structures and Algorithms $\bullet$ Algebraic Structures $\bullet$ Automata and Computability $\bullet$ Computer Systems $\bullet$ Algorithms and Programming $\bullet$ Introduction to AI and Machine Learning $\bullet$ Introduction to Natural Language Processing (PG) $\bullet$ Introduction to Electrical and Electronics Engineering

\section*{Languages}
English $\bullet$ Telugu $\bullet$ Hindi

\end{document}
